% ---------------------------------------------------------------------
% Section 9
% ---------------------------------------------------------------------

\begin{cnread}
电动力学的另一面是束缚态：正负电荷若既有吸引、又有切向运动，就既不湮灭
也不逃逸，而是形成周期轨道。在 \AEG{} 里这正是\textbf{椭圆型}不动点：连续
分数迭代 $z_{k+1}=E-e^2/z_k$ 的判别式 $E^2-4e^2<0$ 时，两个不动点落在
上半平面内部的共轭复点上，算子变成旋转。把这条闭合轨道水平提升到切触空间，
它绕一圈不闭合，而是沿着 Reeb 纤维漂移 $\Delta A=-\oint a\,dM$。把这个漂移
取成基本作用量的整数倍（Boothby--Wang 预量化），就得到玻尔--索末菲条件，
代入库仑势直接给出 $E_n=-\dfrac{m\alpha_e^2}{2\hbar^2 n^2}$。本节末尾说明
这一推导的\textbf{半经典}身份。
\end{cnread}

\subsection{Hyperbolic, parabolic, and elliptic histories}

A projective motion $F\in\PSL_2(\R)$ is classified by its fixed-point set, and by
Paper~0, Subsection~5.5, the finite fixed points solve
\begin{equation}\label{eq:mobius-fixed}
  Cz^{2}+(D-A)z-B=0 .
\end{equation}

\begin{enumerate}[label=(\roman*)]
  \item \emph{Hyperbolic}: two distinct real fixed points.  The motion is a
  translation along the geodesic joining them; in the flow \eqref{eq:theta-flow}
  the scale coordinate advances monotonically.  A pair of charges in radial
  infall or escape follows this type.
  \item \emph{Parabolic}: one double fixed point, necessarily on the boundary.
  The corresponding motion is the horocyclic translation $\mathsf X_s$ of
  \eqref{eq:grid-motions}.
  \item \emph{Elliptic}: two complex conjugate fixed points in
  $\CC\setminus\R$, i.e. inside $\HH^2$ when read in the upper half-plane model.
  The motion is conjugate to a rotation, and its orbit is a closed curve around
  the interior fixed point.  A bound two-body motion follows this type.
\end{enumerate}

The Coulomb problem supplies the elliptic case through the continued-fraction map
of Paper~0, Subsection~5.4.  There, the stationary right-reciprocal update is
$F(z)=a+b/z$ with matrix $M_F=\begin{pmatrix}a&b\\1&0\end{pmatrix}$ and fixed
points $z^{2}-az-b=0$.  Writing the bound-state iteration as
\begin{equation}\label{eq:coulomb-iteration}
  z_{k+1}=E-\frac{e^{2}}{z_k},
  \qquad
  M_F=\begin{pmatrix}E&-e^{2}\\1&0\end{pmatrix},
\end{equation}
the fixed-point equation becomes $z^{2}-Ez+e^{2}=0$, whose discriminant is
$\Delta=E^{2}-4e^{2}$.  For $E<2e$ the discriminant is negative, the two fixed
points are complex conjugates, and the motion is elliptic.

\begin{remark}[Status of the continued fraction]
\label{rem:cf-status}
Equation~\eqref{eq:coulomb-iteration} is the algebraic statement of Paper~0,
Subsection~5.4 with $(a,b)=(E,-e^{2})$.  Its identification with the spectral
problem of a Coulomb Hamiltonian is a dictionary statement, not a theorem: what
\AEG{} supplies is the conjugacy classification of the iteration, and what the
dictionary supplies is the assertion that the physical bound states sit at the
parameter values where that classification becomes elliptic.
\end{remark}

\subsection{The Reeb holonomy of a closed orbit is the action}

Let $\gamma_{\mathrm{orb}}$ be a closed orbit of a bound motion in
$\HH^2_{(M,a)}$, and let $\tilde\gamma_{\mathrm{orb}}$ be its horizontal lift to
$\Cont$.  By Theorem~\ref{thm:dual-tearing}(2) the lift does not close in
$\Cont$: it ends at a Reeb height differing from its starting height by
\begin{equation}\label{eq:orbit-holonomy}
  \Delta A(\gamma_{\mathrm{orb}})
  =-\oint_{\gamma_{\mathrm{orb}}}a\,\mathrm dM
  =-\iint_{\Omega_{\mathrm{orb}}}\mathrm da\wedge\mathrm dM ,
\end{equation}
where $\Omega_{\mathrm{orb}}$ is the region enclosed by the orbit in the
$(M,a)$-plane.

\begin{lemma}[The action slot]
\label{lem:action-slot}
Assume Hypothesis~\ref{hyp:depth} and let the assignment slot of the contact
form be filled by the scale-corrected radial momentum, $a:=r\,p_r$.  Then
\begin{equation}\label{eq:action-form}
  -\oint a\,\mathrm dM=-\oint p_r\,\mathrm dr ,
\end{equation}
so that the Reeb drift \eqref{eq:orbit-holonomy} is minus the classical
mechanical action of the orbit.
\end{lemma}

\begin{proof}
$M=\log(r/\ell_0)$ gives $\mathrm dM=\mathrm dr/r$, hence
$a\,\mathrm dM=r p_r\,\mathrm dr/r=p_r\,\mathrm dr$; the sign is the orientation
of the loop.
\end{proof}

\begin{warning}[What the substitution assumes]
\label{warn:action-slot}
Lemma~\ref{lem:action-slot} is the point at which the note departs from purely
\AEG{} statements: the assignment of the contact form is a number attached to a
history, and calling it $r p_r$ --- with $r$ the physical separation --- is a
dictionary step additional to Hypotheses~\ref{hyp:charge}--\ref{hyp:pairing}.  It
is however the least avoidable one, because without it the Reeb drift is a
dimensionless area and cannot be compared with $\hbar$.
\end{warning}

\subsection{Monodromy and the Bohr--Sommerfeld condition}

Suppose the Reeb fiber is compactified.  By the Boothby--Wang theorem
\cite{BoothbyWang1958,Geiges2008}, if the contact form's Reeb field is a
circle action with period $2\pi\hbar$, then $\Cont$ is the total space of a
principal $S^{1}$-bundle --- the prequantum bundle of the underlying symplectic
manifold $\HH^2$ --- and the horizontal lift of a closed orbit closes in the
fiber exactly when its drift is an integer multiple of the period.

\begin{definition}[Monodromy condition]
\label{def:monodromy}
The quantized contact model is the quotient $\Cont/(2\pi\hbar\,\mathbb Z\cdot R)$
of the Reeb flow.  A closed orbit $\gamma_{\mathrm{orb}}$ defines a monodromy
element $\exp(\mathrm i\,\Delta A(\gamma_{\mathrm{orb}})/\hbar)\in S^{1}$, and it
defines a single-valued state exactly when
\begin{equation}\label{eq:bohr-sommerfeld}
  \oint_{\gamma_{\mathrm{orb}}}(-a\,\mathrm dM)
  =\oint p_r\,\mathrm dr
  =2\pi n\hbar ,
  \qquad n\in\mathbb Z .
\end{equation}
\end{definition}

\begin{theorem}[Coulomb spectrum]
\label{thm:spectrum}
Under Hypothesis~\ref{hyp:charge}--\ref{hyp:pairing}, Lemma~\ref{lem:action-slot},
and Definition~\ref{def:monodromy}, the bound orbits of a Coulomb pair in the
radial coordinate satisfy
\begin{equation}\label{eq:radial-action}
  \oint p_r\,\mathrm dr
  =2\pi\Bigl(\frac{m\alpha_e}{\sqrt{-2mE}}-l\Bigr)=2\pi n_r\hbar ,
\end{equation}
with $\alpha_e=e^{2}/(4\pi)$ and $l$ the angular momentum.  Writing
$n=n_r+l/\hbar$, the energies are
\begin{equation}\label{eq:spectrum}
  E_n=-\frac{m\alpha_e^{2}}{2\hbar^{2}n^{2}},
  \qquad n=1,2,3,\dots
\end{equation}
\end{theorem}

\begin{proof}
The second equality in \eqref{eq:radial-action} is the standard Coulomb radial
action; Appendix~\ref{app:action} evaluates the half-action
$\int_{r_-}^{r_+}p_r\,\mathrm dr=\pi\bigl(m\alpha_e/\sqrt{-2mE}-l\bigr)$ both
in closed form and by numerical quadrature at three parameter points
(Checks~V7.1).  Dividing by $2\pi\hbar$ and putting $n_r=n-l/\hbar$ gives
$m\alpha_e/\sqrt{-2mE}=\hbar n$, hence $-2mE=m^{2}\alpha_e^{2}/(\hbar^{2}n^{2})$
and \eqref{eq:spectrum}.  Check~V7.2 verifies the quantization condition and
Check~V7.3 the ground-state value.
\end{proof}

\begin{remark}[The residue picture]
\label{rem:residue}
The integral \eqref{eq:radial-action} is a contour integral between the two
turning points, i.e. between the two real roots of the quadratic that Paper~0,
Proposition~5.4 attaches to the iteration.  Its closed form is the residue
difference between the two ends of the contour.  In the bivaluation language of
Section~\ref{sec:radiation} the two ends are the two punctures
$(q^+,q^-)=(0,\infty)$ of the reference pair, so the same integral is a residue
pairing at the two ends of the frame.  This is a suggestive identification, not a
proof; it is recorded because it is the natural continuation of the present
picture.
\end{remark}

\subsection{How semiclassical is this?}

\begin{warning}[Range of the derivation]
\label{warn:semiclassical}
Theorem~\ref{thm:spectrum} is a Bohr--Sommerfeld (EBK) computation, not a
solution of the Schr\"odinger equation.  Three limitations should be stated.
\begin{enumerate}[label=(\alph*)]
  \item The condition \eqref{eq:bohr-sommerfeld} quantizes a single action
  integral; a genuine quantization would need the full polarization and the
  half-form (metaplectic) correction of geometric quantization.  For the Coulomb
  problem those corrections are known to conspire so that the leading
  Bohr--Sommerfeld spectrum is exact, but that is a property of this particular
  potential, not a consequence of the present construction.
  \item The Langer-type correction $l\mapsto l+\tfrac12$ changes the spectrum
  beyond leading order; no statement about it is made here.
  \item Nothing above defines a Hilbert space, an inner product, or an operator
  algebra.  The Reeb circle of Definition~\ref{def:monodromy} is a compactification
  hypothesis, and the multiplicities of the resulting levels are not counted.
\end{enumerate}
What the section does establish is narrower and exact: the quantity that
Bohr--Sommerfeld quantizes is, in this picture, the Reeb drift of a closed
horizontal loop --- an object that exists in \AEG{} independently of any physical
reading --- and the Coulomb input enters only through
Lemma~\ref{lem:action-slot} and the classical integral of
Theorem~\ref{thm:spectrum}.
\end{warning}
